\section{Related Work}
\label{sec:related}

\paragraph{Controlled visual task generation.}
CLEVR established synthetic scene graphs and executable question programs as a
way to control compositional visual reasoning; GQA extended structured
programs to questions over real-image scene graphs
\citep{johnson2017clevr,hudson2019gqa}. PuzzleVQA procedurally generates
abstract patterns for diagnosing perceptual, inductive, and deductive failures
\citep{chia2024puzzlevqa}. Task Me Anything instead organizes programmatic
multimodal evaluation around an extensible taxonomy of assets, attributes,
relations, and task plans \citep{zhang2024taskmeanything}. WorldBench takes a
real-image route, using a broad concept taxonomy to curate questions that
require knowledge and reasoning beyond recognition \citep{yin2026worldbench}.
These systems demonstrate the value of explicit structure for generation and
coverage analysis; \trace additionally makes the executable task contract the
stable unit of generation, sampling, verification, and analysis.

\paragraph{Executable generator--verifier environments.}
Reasoning Gym and Enigmata show how programmatic generators and exact
verifiers can produce scalable reasoning curricula
\citep{stojanovski2025reasoninggym,chen2025enigmata}. OpenThoughts likewise
emphasizes systematic choices about sources, filtering, mixing, and teacher
generation rather than treating a training corpus as an opaque collection
\citep{guha2025openthoughts}. The former systems center executable textual
reasoning, whereas the latter emphasizes transparent corpus construction and
mixture design. Together they motivate training environments whose data and
verification decisions remain inspectable rather than being fixed only in a
static corpus. \trace brings this generator--verifier model to rendered
visual scenes while retaining stable task identities across regenerated data.

\paragraph{Open data for multimodal reinforcement learning.}
Recent visual RL systems pair heterogeneous training collections with
rule-based or task-specific rewards. Visual-RFT introduced rewards for visual
classification, detection, and spatial tasks; Vision-R1 and LMM-R1 constructed
cold-start and RL data for multimodal reasoning; and MM-Eureka, R1-Onevision,
Reason-RFT, and VLM-R1 cover visual mathematics, counting, structure
perception, spatial transformation, and referring expressions
\citep{liu2025visualrft,huang2025visionr1,peng2025lmmr1,
meng2025mmeureka,yang2025r1onevision,tan2025reasonrft,shen2025vlmr1}.
Vero scales this direction to a 600K-example mixture from 59 datasets with
task-routed rewards and controlled mixture analysis \citep{sarch2026vero}.
Collectively, these works establish that broad multimodal RL can benefit from
specialized rewards and carefully designed mixtures. Their task identities and
granularity are generally inherited from source datasets, however, so a source
weight need not correspond to a uniform reasoning unit. \trace instead builds
the mixture from explicit task programs and balances sampling over objective
contracts.

\paragraph{Mixture selection and curricula.}
Beyond constructing RL examples, several systems study which examples to use
and in what order. Curr-ReFT combines difficulty-aware progression with
rejected-sampling self-improvement, while SoTA with Less and RAP select compact
or high-value subsets for multimodal post-training
\citep{deng2025currreft,wang2025sotawithless,li2025rap}. MoDoMoDo formulates
multidomain RLVR as a data-mixture optimization problem
\citep{liang2025modomodo}, and Vision-G1 combines influence-based filtering
across 46 sources with a multi-round curriculum \citep{zha2026visiong1}.
WeThink and DeepVision-103K further broaden RL-ready data through scalable
synthesis, curation, verification, and difficulty stratification
\citep{yang2025wethink,sun2026deepvision}. These approaches show that coverage
alone is insufficient: source interaction, difficulty, and example quality
also shape learning. \trace is complementary in providing stable task
contracts as a finer-grained unit for such sampling and curriculum methods.

\paragraph{Synthetic and verifiable visual RL data.}
The closest systems construct new visual experience specifically for
reinforcement learning. Jigsaw-R1 uses controlled jigsaw puzzles to study
rule-based visual RL, while VisualSphinx expands seed puzzle rules into a
large rule-to-image corpus of visual logic problems
\citep{wang2025jigsawr1,feng2025visualsphinx}. Game-RL adapts executable game
code into verifiable multimodal tasks, and COGS decomposes and recomposes
perception and reasoning factors for charts and webpages
\citep{tong2025gamerl,gu2025cogs}. SynthRL produces verified harder variants
of visual-mathematics questions, whereas ViCrit injects localized,
exactly-checkable errors into image descriptions as a visual-perception proxy
task \citep{wu2025synthrl,wang2025vicrit}. Sphinx procedurally generates 25
visual reasoning task families with deterministic answers and demonstrates
transfer from RLVR training \citep{alam2026sphinx}. Concurrent TRON is the
closest environment-level system: it provides 520 online generator--verifier
environments in five ability buckets, with difficulty-controlled rollouts and
audits of generation reliability, diversity, and difficulty
\citep{yang2026tron}.

\begin{table}[h]
  \centering
  \scriptsize
  \setlength{\tabcolsep}{2.5pt}
  \renewcommand{\arraystretch}{1.08}
  \caption{Reported structural scale of procedural visual RL systems. Counts
  follow each work's authored units; generated examples and LLM-expanded rules
  are not counted as objectives.}
  \label{tab:environment-comparison}
  \begin{tabular}{@{}l l c l r l@{}}
    \toprule
    System & Scope & Groups & Visual units & Objectives & Relation \\
    \midrule
    Jigsaw-R1    & Jigsaw       & 1  & 1 grammar    & 2 forms & joint \\
    VisualSphinx & Logic puzzle & 8  & --           & 1       & joint \\
    GameQA       & Games        & 4  & 30 games     & 158     & game$\to$task \\
    Sphinx       & Abstract     & 5  & 25 renderers & 25      & 1:1 \\
    TRON         & Multiskill   & 5  & 520 envs     & 520     & joint \\
    \textbf{\trace} & Multidomain & \textbf{11} & \textbf{277 scenes}
      & \textbf{1,000} & scene$\to$task \\
    \bottomrule
  \end{tabular}
\end{table}

VisualSphinx's 41,287 retained LLM-generated rules instantiate one completion
objective, and its RL run samples 10,000 examples. \trace is distinguished by
both its 11-domain breadth and its explicit 277-scene/1,000-task split. GameQA
also separates games from tasks within one domain, while TRON couples generation
and verification in 520 joint environments; COGS and SynthRL do not report
fixed task inventories.
